\section{Síntesis de los resultados principales}
\label{sec:sintesis_conclusiones}

Este trabajo ha evaluado la aplicación de \emph{Differential Machine Learning} (DML) a la valoración y cobertura de una opción asiática aritmética sobre capacidad de cómputo, utilizando como exposición económica el coste por GPU-hora. El diseño combina datos observados del mercado H100, una dinámica Log-OU, simulación RQMC y una comparación homogénea entre un MLP convencional y un modelo DML con la misma arquitectura y protocolo de entrenamiento. El objetivo principal ha sido determinar si la información diferencial mejora la aproximación del precio y, especialmente, de la Delta, y si esa mejora se traduce en una reducción efectiva del riesgo de cobertura.

La serie H100 analizada presenta un comportamiento compatible con reversión a la media durante el periodo muestral. La estimación física del Log-OU proporciona $\widehat{\kappa}_P=32{,}6645$, equivalente a una vida media aproximada de $7{,}75$ días. El bootstrap muestra, sin embargo, una incertidumbre relevante en torno a este parámetro. Por ello, el Log-OU debe interpretarse como una representación parsimoniosa y operativa del entorno de simulación, no como la dinámica verdadera del mercado de compute.

En \textbf{pricing}, DML no domina uniformemente al MLP. Con $65\,536$ estados de entrenamiento, DML obtiene un RMSE de test de $0{,}001516$ frente a $0{,}001624$ del MLP, mientras que el MLP presenta un MAE ligeramente inferior. En el conjunto de referencia, construido con etiquetas RQMC de mayor precisión, DML obtiene una ventaja más clara, especialmente en RMSE. La conclusión adecuada es, por tanto, que DML mantiene una precisión de precio comparable y mejora algunas métricas, pero no ofrece una superioridad general en valoración.

La evidencia es mucho más concluyente para la \textbf{Delta}. En la configuración final, el RMSE disminuye de $0{,}008528$ con MLP a $0{,}001817$ con DML y el MAE de $0{,}001234$ a $0{,}000216$. Esto equivale a reducciones emparejadas medias próximas al $78{,}5\%$ y al $82{,}0\%$, respectivamente. DML obtiene además menor MAE y RMSE de Delta en las cinco semillas y en todos los tamaños de entrenamiento considerados, patrón que también se reproduce sobre el conjunto de referencia. La principal evidencia empírica del trabajo es, por tanto, una mejora amplia, consistente y robusta de la sensibilidad cuando se incorpora supervisión diferencial.

Los resultados muestran también una mayor \textbf{eficiencia muestral} dentro de la malla experimental: DML entrenado con $1\,024$ estados supera al MLP con $16\,384$ en MAE y RMSE de Delta. Este resultado no permite establecer una ratio universal de eficiencia, pero sí indica que las derivadas aportan información adicional útil por estado simulado. Esta ventaja no implica menor coste computacional total, ya que el entrenamiento diferencial es más costoso que el convencional.

El análisis de errores extremos refuerza esta conclusión y, al mismo tiempo, identifica una debilidad concreta. Los mayores errores de Delta se concentran desproporcionadamente cerca de las fronteras mensuales de fixing: la distancia mediana al fixing más próximo es de aproximadamente $0{,}49$ días entre los cincuenta mayores errores de DML, frente a $7{,}54$ días en la muestra completa. Esta evidencia sugiere que la representación temporal utilizada captura de forma imperfecta los cambios discretos del estado contractual.

La utilidad económica de la Delta aprendida se confirma claramente en la \textbf{cobertura local}. Para un shock del $1\%$, el RMSE del residuo pasa de $6{,}6878\times10^{-6}$ con MLP a $1{,}0584\times10^{-6}$ con DML, una mejora aproximada del $84{,}2\%$; frente a la posición sin cobertura, la reducción supera el $99\%$. El mismo orden se mantiene para shocks del $0{,}5\%$ y del $2\%$. En este entorno, una mejor estimación de la sensibilidad se traduce directamente en una mejor neutralización del riesgo de primer orden.

La \textbf{cobertura dinámica} introduce una conclusión más exigente. DML sigue siendo muy superior al MLP, con RMSE terminal de $0{,}038698$ frente a $0{,}106515$, y obtiene menor error absoluto que el MLP en el $70{,}07\%$ de las trayectorias. También mejora frente a no cubrir en el $77{,}33\%$ de las trayectorias y reduce notablemente la mediana del error absoluto. Sin embargo, su RMSE permanece por encima del $0{,}014998$ de la posición sin cobertura. El benchmark Oracle, con RMSE próximo a $0{,}0023$, confirma que el esquema de hedge puede reducir riesgo dentro del modelo; el deterioro de los modelos aprendidos procede de una pequeña fracción de trayectorias extremas que domina el error cuadrático agregado.

En conjunto, el resultado central del TFM es que \textbf{DML mejora de forma sustancial la Delta y esa mejora posee valor económico en cobertura local, pero una Greek más precisa no garantiza por sí sola una mejor cobertura dinámica}. La eficacia final depende también del instrumento utilizado, de su sensibilidad al spot, del rebalanceo discreto, del estado contractual y del comportamiento de las colas.

\section{Respuesta a las preguntas de investigación}
\label{sec:respuesta_preguntas}

La primera pregunta recibe una respuesta diferenciada: la supervisión de Delta
mejora de forma consistente la sensibilidad, mientras que su efecto sobre el precio
depende de la métrica y del conjunto evaluado. La segunda recibe una respuesta
condicionada por el uso: la ventaja se transmite al residuo local, pero no basta
para superar a no cubrir en RMSE terminal con la política dinámica ejecutada.
Estos resultados delimitan cuándo resulta útil añadir información diferencial.

\section{Implicaciones financieras y aportación}
\label{sec:implicaciones_financieras}

Los resultados sugieren que la evaluación de un surrogate para derivados no debería limitarse al error de precio cuando su uso posterior depende de las Greeks. Dos modelos pueden presentar errores de valoración similares y, sin embargo, ofrecer una calidad muy distinta para gestión del riesgo. En este TFM, esa diferencia es precisamente la principal ventaja de DML: la incorporación de información diferencial apenas altera la conclusión sobre pricing, pero mejora de forma sustancial la sensibilidad relevante para cobertura.

La segunda implicación es que la eficiencia muestral debe distinguirse de la eficiencia computacional. Las etiquetas diferenciales permiten extraer más estructura de cada estado simulado y alcanzar una mejor precisión de Delta con menos observaciones dentro del experimento realizado. Sin embargo, el entrenamiento DML requiere un coste adicional. La utilidad del método dependerá, por tanto, de que las derivadas puedan obtenerse de forma fiable y de que las sensibilidades tengan suficiente importancia económica para justificar dicho coste.

La comparación entre cobertura local y dinámica aporta una tercera implicación. Una mejor Delta constituye una condición favorable para el hedge, pero no suficiente. En la estrategia dinámica la sensibilidad debe transformarse en unidades de un forward modelizado mediante su Jacobiano respecto al spot; bajo fuerte reversión a la media, un Jacobiano reducido puede amplificar errores residuales de Delta. El diseño combina además rebalanceo discreto y transiciones de fixing, sin identificar contribuciones causales separadas ni un factor de base independiente. Por este motivo, la calidad de una Greek debe evaluarse conjuntamente con el mecanismo mediante el que se convierte en una posición económica.

Finalmente, el trabajo muestra que el coste por GPU-hora puede formularse como una exposición cuantitativa susceptible de modelización mediante un payoff dependiente del promedio. Esto no implica que exista actualmente un mercado líquido y estandarizado de derivados sobre compute ni que los precios obtenidos constituyan cotizaciones implementables. La aportación es metodológica: integrar datos de compute, una dinámica estocástica, valoración RQMC, aprendizaje de precio y Delta, y evaluación de cobertura dentro de un protocolo común que permite identificar con precisión dónde DML aporta valor y dónde esa ventaja deja de ser suficiente.

\section{Limitaciones}
\label{sec:limitaciones}

La principal limitación procede de los \textbf{datos}. La serie H100 utilizada contiene $92$ observaciones diarias entre el 26 de mayo y el 25 de agosto de 2026. Esta historia es suficiente para una primera caracterización, pero corta para establecer la estabilidad de largo plazo de la reversión a la media. El bootstrap confirma una incertidumbre apreciable en $\kappa_P$. Además, las referencias de Ornn y SMM no corresponden a contratos idénticos y sólo permiten una comparación descriptiva entre mercados. H100 tampoco representa de forma exhaustiva el conjunto del mercado de capacidad computacional.

La segunda limitación afecta a la \textbf{especificación del subyacente}. El Log-OU de un factor proporciona positividad, reversión a la media y transición exacta, pero excluye volatilidad estocástica, saltos, estacionalidad, factores tecnológicos o regionales y posibles cambios de régimen. En consecuencia, el experimento identifica con precisión el error de aproximación de MLP y DML respecto a la función definida por el modelo, pero no cuantifica el error de especificación que existiría frente a un futuro mercado real.

Una tercera limitación es la \textbf{identificación de la medida de valoración}. Al no existir una historia suficientemente líquida de forwards y opciones sobre GPU-hour que permita inferir una medida neutral al riesgo de mercado, los parámetros utilizados bajo $\mathbb Q$ son una especificación modelizada. Las bandas empleadas no deben interpretarse como intervalos de confianza de parámetros neutrales al riesgo y los precios sintéticos no equivalen a cotizaciones observables. Esta restricción afecta al nivel absoluto de valoración, aunque menos a la comparación interna entre MLP y DML, que utilizan la misma ley de pricing y las mismas etiquetas.

La cuarta limitación reside en la \textbf{representación del estado contractual}. El vector de entrada incorpora el tiempo a vencimiento y el promedio acumulado, pero no variables explícitas como el número de fixings pendientes o la distancia al siguiente fixing. Con el calendario fijo, $\tau$ determina ya esos metadatos: explicitarlos facilita la representación, sin añadir información. La concentración de errores extremos cerca de estas fronteras sugiere que una red suave debe aprender implícitamente cambios contractuales discretos a partir de una variable temporal continua. El diagnóstico fue realizado después de congelar los modelos y no se utilizó para reentrenarlos, preservando la validez de la evaluación final.

La salida neuronal tampoco garantiza precio no negativo ni monotonicidad respecto al spot. Las violaciones observadas se documentan sin recorte retrospectivo; una arquitectura restringida deberá contrastarse con un nuevo test. Además, las pruebas local y dinámica cambian la distribución temporal evaluada (estados interiores frente a fronteras), por lo que su comparación no aísla el efecto del rebalanceo. La muestra dinámica es condicional a permanecer en el dominio y la dispersión entre semillas sólo recoge variación del entrenamiento sobre los mismos datos.

La quinta limitación afecta al \textbf{instrumento de cobertura}. El forward de compute utilizado es modelizado y no incorpora directamente bid--ask, liquidez, costes de ejecución, margen o límites de posición. Además, su sensibilidad al spot puede ser pequeña bajo reversión a la media, amplificando errores de Delta al convertirlos en unidades de hedge. La corrección de descuento del ratio forward se presenta como diagnóstico posterior y no altera la jerarquía de RMSE. Por ello, los resultados dinámicos deben interpretarse como evidencia dentro del entorno modelizado y no como estimación directa del rendimiento de una estrategia negociable.

En síntesis, la comparación MLP--DML posee una validez interna elevada porque ambos modelos comparten datos, función de valoración y protocolo. La validez externa del nivel de precios y de la cobertura dinámica es más limitada y depende de la representatividad de los datos, de la especificación de $\mathbb Q$ y de la existencia futura de instrumentos observables sobre compute.

\section{Líneas futuras de investigación}
\label{sec:future_work}

Las extensiones prioritarias deben responder directamente a las limitaciones identificadas. La primera es enriquecer el estado contractual con variables explícitas como $N_{\mathrm{future}}$ y $d_{\mathrm{fix}}$ y comprobar, sobre un nuevo conjunto de test, si reducen la concentración de errores alrededor de los fixings. La segunda es reestimar la dinámica a medida que se disponga de una historia H100 más extensa y evaluar la estabilidad temporal de $\kappa$, $\theta$ y $\sigma$ antes de introducir modelos más complejos.

Una tercera extensión consiste en sustituir progresivamente el forward modelizado por instrumentos observables cuando exista suficiente liquidez, incorporando precios de mercado, costes de transacción y riesgo de base real. Esos instrumentos permitirían además identificar mejor la dinámica de valoración y contrastar precios modelizados con cotizaciones efectivas.

Sólo después de disponer de una base empírica más rica tendría sentido ampliar el Log-OU a especificaciones multifactoriales, estudiar otras GPUs o regiones y comparar el enfoque actual con estrategias de \emph{Deep Hedging} que aprendan directamente la política de cobertura. La prioridad no es añadir complejidad algorítmica, sino resolver primero las limitaciones de datos, estado e instrumentación que los propios resultados han puesto de manifiesto.
